\documentclass[12pt,a4paper]{article}
\usepackage[T2A]{fontenc}
\usepackage[utf8]{inputenc}
\usepackage[russian]{babel}
\usepackage{amsmath}
\usepackage{booktabs}
\usepackage{geometry}
\geometry{margin=2cm}

\begin{document}
\begin{center}
    {\LARGE Лабораторная работа 3}
\end{center}
\begin{flushright}
\textbf{Работу выполнили:} \\
Шубин Егор\\
Санатин Дмитрий\\
\textbf{Преподаватель:} \\
Танченко Юлия Валерьевна
\end{flushright}



\section*{Условие}
Проверить гипотезу
\[
H_0: X \sim \mathrm{Exp}(\lambda), \quad \lambda = 1.55
\]
для заданной выборки объема $n=100$ с использованием критерия Пирсона (хи-квадрат),
уровень значимости $\alpha = 0.05$. Использовать равномерную сетку из 10 интервалов.

\section*{Построение интервалов и вероятностей}
Максимум выборки: $x_{\max}=2.03$.  
Берем шаг сетки:
\[
h = \frac{x_{\max}}{10} = 0.203.
\]
Интервалы:
\[
I_i=[ih, (i+1)h),\ i=0,\dots,8, \qquad I_9=[9h, +\infty).
\]

Для экспоненциального распределения:
\[
F(x)=1-e^{-\lambda x}, \quad x\ge 0.
\]
Тогда теоретические вероятности:
\[
p_i=P_{H_0}(X\in I_i)=
\begin{cases}
e^{-\lambda ih}-e^{-\lambda (i+1)h}, & i=0,\dots,8,\\
e^{-\lambda \cdot 9h}, & i=9.
\end{cases}
\]

\section*{Таблица расчетов}
\begin{center}
\small
\begin{tabular}{ccccc}
\toprule
Интервал $I_i$ & Наблюд. $\nu_i$ & Вычисление $p_i$ & Теор. $p_i$ & Ожид. $n p_i$ \\
\midrule
$[0.000;0.203)$ & 34 & $e^{-1.55\cdot 0.000}-e^{-1.55\cdot 0.203}$ & 0.269956 & 26.9956 \\
$[0.203;0.406)$ & 24 & $e^{-1.55\cdot 0.203}-e^{-1.55\cdot 0.406}$ & 0.197080 & 19.7080 \\
$[0.406;0.609)$ & 8  & $e^{-1.55\cdot 0.406}-e^{-1.55\cdot 0.609}$ & 0.143877 & 14.3877 \\
$[0.609;0.812)$ & 11 & $e^{-1.55\cdot 0.609}-e^{-1.55\cdot 0.812}$ & 0.105036 & 10.5036 \\
$[0.812;1.015)$ & 4  & $e^{-1.55\cdot 0.812}-e^{-1.55\cdot 1.015}$ & 0.076681 & 7.6681 \\
$[1.015;1.218)$ & 8  & $e^{-1.55\cdot 1.015}-e^{-1.55\cdot 1.218}$ & 0.055981 & 5.5981 \\
$[1.218;1.421)$ & 7  & $e^{-1.55\cdot 1.218}-e^{-1.55\cdot 1.421}$ & 0.040868 & 4.0868 \\
$[1.421;1.624)$ & 1  & $e^{-1.55\cdot 1.421}-e^{-1.55\cdot 1.624}$ & 0.029836 & 2.9836 \\
$[1.624;1.827)$ & 2  & $e^{-1.55\cdot 1.624}-e^{-1.55\cdot 1.827}$ & 0.021781 & 2.1781 \\
$[1.827;+\infty)$ & 1 & $e^{-1.55\cdot 1.827}$ & 0.058904 & 5.8904 \\
\bottomrule
\end{tabular}
\normalsize
\end{center}

\section*{Статистика критерия}
\[
\eta_{\text{набл}}=\sum_{i=1}^{10}\frac{(\nu_i-np_i)^2}{np_i}=15.8668.
\]

Число степеней свободы:
\[
r = k-1-s = 10-1-0=9.
\]

\[
\delta(X)=
\begin{cases}
H_0, & \eta \le c,\\
H_1, & \eta > c.
\end{cases}
\]
\[
c=\chi^2_{1-\alpha;\,r},
\]

\[
\chi^2_{1-\alpha;\,r}=\chi^2_{0.95;\,9}=16.919.
\]

Критическая область правосторонняя:
\[
\mathcal{D}=\left\{\eta>\chi^2_{1-\alpha;\,r}\right\}.
\]
Сравнение:
\[
\eta_{\text{набл}}=15.8668 < 16.919=\chi^2_{0.95;\,9}.
\]

\section*{Вывод}
На уровне значимости $\alpha=0.05$ нет оснований отвергнуть гипотезу $H_0$.
Выборка согласуется с экспоненциальным распределением с параметром $\lambda=1.55$.

\end{document}
